\documentclass{amsart}
\usepackage{amssymb}
\usepackage{hyperref}
\usepackage{amsrefs}

\newcommand{\Q}{\mathbb{Q}}
\newcommand{\QQ}{\mathbb{Q}}
\newcommand{\Qp}{\mathbb{Q}_p}
\newcommand{\Z}{\mathbb{Z}}
\newcommand{\ZZ}{\mathbb{Z}}
\newcommand{\Zp}{\mathbb{Z}_p}
\newcommand{\hatZ}{\hat{\Z}}

\DeclareMathOperator{\Aut}{Aut}
\DeclareMathOperator{\Gal}{Gal}

\newcommand{\GQp}{\Gal(\overline{\Q}_p/\Qp)}

\newtheorem*{theorem}{Theorem}

%Problem statement
%Solution verification script
%Description of importance to the field
%Qualitative assessment of difficulty

\title{Absolute Galois group of $\Q_2$}
\author{David Roe}

\begin{document}
\begin{abstract}
Find a presentation for the absolute Galois group of the p-adic field $\Q_2$ as a profinite group.
\end{abstract}

\section{Problem statement}

Give a presentation for the absolute Galois group of the field of $2$-adic numbers as a profinite group.

\section{Overview}

The absolute Galois group $\Gal(\overline{K}/K)$ of a field $K$ is the projective limit of all of the finite Galois groups $\Gal(E/K)$.  It packages together the information about all finite extensions, and studying the Galois group of the rational field $\mathbb{Q}$ is a central problem in algebraic number theory.  One method for approaching this group is to study the analogous Galois groups of p-adic fields $\mathbb{Q}_p$.  In this case, there is an explicit presentation of $\GQp$ for $p>2$, described in Section \ref{sec:background}.  For $p=2$, we have a description for the absolute Galois group of some extensions of $\mathbb{Q}_2$, but not for $\mathbb{Q}_2$ itself.  Finding such a description would fill in a gap in our explicit understanding of Galois groups, and would have consequences in terms of counting p-adic fields with a given Galois group.

\section{Background}
\label{sec:background}

We describe the presentation for known $p > 2$.  Let $h$ be a $(p-1)$st root of unity in $\Zp$.
Let $\pi = \pi_p$ be the element of $\hatZ = \prod_\ell \ZZ_\ell$ with coordinate $1$ in the $\Zp$-component and $0$ in the $\ZZ_\ell$ components for $\ell \ne p$.  Then for $x, y$ in a profinite group\footnote{See \cite{ribes-zalesskii:profinite_groups}, especially sections 3.3 and 4.1, for relevant background on profinite groups.}, set
\[
\langle x, y \rangle = (x^{h^{p-1}}y x^{h^{p-2}} y \cdots x^h y)^{\frac{\pi}{p-1}}.
\]
Write $[x,y]$ for the commutator $xyx^{-1}y^{-1}$.

\begin{theorem}[{\cite{neukirch-schmidt-wingberg:CohomologyOfNumberFields}*{Thm. 7.5.14}}] \label{thm:GQp_desc}
Let $p \ne 2$.  The absolute Galois group $\GQp)$ is isomorphic to the profinite group with $4$ generators $\sigma, \tau, x_0, x_1$, subject to the following conditions and relations.
\begin{enumerate}
\item The closed subgroup topologically generated by $x_0$ and $x_1$ is normal and pro-$p$.
\item The elements $\sigma, \tau$ satisfy the tame relation
\[
\tau^\sigma = \tau^q.
\]
\item The generators satisfy the wild relation
\[
x_0^\sigma = \langle x_0, \tau \rangle x_1^p [x_1, y_1],
\]
where $y_1$ is defined below.
\end{enumerate}

In order to define $y_1$, let $\Qp^t$ be the maximal tamely ramified extension of $\Qp$ and define
$\beta : \Gal(\Qp^t/\Qp) \to \Zp^\times$ by setting $\beta(\sigma) = 1$ and $\beta(\tau) = h$.
For $\rho$ in the subgroup generated by $\sigma$ and $\tau$ and an arbitrary $x$, set
\[
\{x, \rho\} = (x\rho^2x^{\beta(\rho)}\rho^2\dots x^{\beta(\rho^{p-2})}\rho^2)^{\frac{\pi}{p-1}}.
\]
Similar to $\pi_p$ above, let $\pi_2 \in \hatZ$ have a $1$ in the $\Z_2$-component and $0$ in every other component, and set $\tau_2 = \tau^{\pi_2}$ and $\sigma_2 = \sigma^{\pi_2}$.
Set
\[
y_1 = x_1^{\tau_2^{p+1}}\{x_1, \tau_2^{p+1}\}^{\sigma_2\tau_2^{(p-1)/2}}\{\{x_1, \tau_2^{p+1}\}, \sigma_2\tau_2^{(p-1)/2}\}^{\sigma_2\tau_2^{(p+1)/2}+\tau_2^{(p+1)/2}}.
\]
\end{theorem}

\section{Importance}

Galois representations are at the core of many problems in modern number theory.  There are many perspectives on Galois representations, but if you have access to a presentation of an absolute Galois group then they can be viewed very concretely.  The case of $p=2$ often needs to be handled separately in number theory, and global applications often require the ability to handle all primes.

\section{Difficulty}

The last major work on finding presentations for $p$-adic fields was done in the 1980s.  To a large extent, I believe that the difficulty is that the answer is likely to be messy, as evidenced by the complicated relation needed in the case that $p$ is odd.  Hopefully, this feature makes the problem a good target for solution by an AI since I don't expect major conceptional developments to be required.

\section{Submitting a presentation}

A submission for this problem consists of a file with the following form:

\vspace{0.1in}
\texttt{variables}
\vspace{0.1in}

\texttt{definitions}
\vspace{0.1in}

\texttt{relations}
\vspace{0.1in}

\begin{itemize}
\item Sections should be separated by single blank lines, and the \texttt{definitions} section is optional. 
\item The \texttt{variables} should consist of a comma separated list of variable names on one line corresponding to the generators of the presentation.  The first and second generators (say $\sigma$ and $\tau$ respectively) are distinguished, and must satisfy the implicit relation $\tau^\sigma = \tau^p$ (which should not be included in the \texttt{relations} section).  All later generators should be contained within the $p$-core of the group (the intersection of all $p$-Sylow subgroups).
\item In the \texttt{definitions} section, you may define additional variables in terms of the generators and variables defined earlier within the \texttt{definitions} section.  Each line should take the form \texttt{var = expr}, where \texttt{expr} is an expression built using group operations.
\item In the \texttt{relations} section, you should give the relations as expressions in terms of variables defined in the previous two sections.  If there are multiple relations, give one on each line.
\end{itemize}

Expressions in profinite groups may involve exponents lying in $\hat{\Z}$.  To specify such a constant, you should provide its reduction modulo $85667662080 = 2^{8} \cdot 3^{2} \cdot 5 \cdot 7 \cdot 11 \cdot 13 \cdot 17 \cdot 19 \cdot 23$, which is the least common multiple of the exponents of all the groups being used to verify the presentation.

\section{Verification}

Verification of the proposed presentation consists of using it to count extensions of $\Q_2$ with specified finite Galois group using methods described in \cite{roe:19a}.  These counts are then compared to known counts from two sources:
\begin{itemize}
\item The L-functions and modular forms database \cite{LMFDB}, which contains all $2$-adic fields of degree up to 23 together with their Galois groups.
\item The known maximal pro-$2$ quotient, which allows for counting the number of fields with Galois group any $2$-group.  The counts are also computed using the techniques in \cite{roe:19a}, though it is also possible to use the methods of \cite{yamagishi:95a} that generalize Shafaravich's formula \cite{shafarevich:47a} for odd $p$ counting the number of extensions of $\Qp$ with Galois group $G$ (requiring $d$ generators):
\[
\frac{1}{\lvert \Aut(G) \rvert} \left(\frac{\lvert G \rvert}{p^d}\right)^{n+1} \prod_{i=0}^{d-1} (p^{n+1} - p^i).
\]
\end{itemize}

Since we only check finitely many groups, this verification is not a proof that the submission is correct, but it will provide good evidence that a solution developed using a theoretical approach does not have a basic flaw.  In particular, agents should not have direct access to the set of groups used for verification when constructing their solution.

\begin{thebibliography}{9}

\bibitem{iwasawa:55a} Kenkichi Iwasawa, \textit{On galois groups of local fields}, Transactions of the AMS \textbf{80} (1955), 448--469.
\bibitem{LMFDB} The LMFDB Collaboration. ``The L-functions and Modular Forms
Database''. \url{https://www.lmfdb.org}.  Online, accessed Jan 26, 2026.
\bibitem{neukirch-schmidt-wingberg:CohomologyOfNumberFields} J\"urgen Neukirch, Alexander Schmidt, and Kay Wingberg, \textit{Cohomology of number fields}, 2nd ed., Springer-Verlag, Berlin, 2015.
\bibitem{ribes-zalesskii:profinite_groups} Luis Ribes and Pavel Zalesskii, \textit{Profinite groups}, 2nd ed., Springer-Verlag, Berlin, 2010.
\bibitem{roe:19a} David Roe, \textit{The Inverse Galois Problem for $p$-adic fields}, Proceedings of the Thirteenth Algorithmic Number Theory Symposium (ANTS-XIII), Open Book Series 2, Math. Sci. Pub., Berkeley, 2019, pp 393-409.
\bibitem{shafarevich:47a} Igor Shafarevich, \textit{On p-extensions}, Mat. Sb. \textbf{20} (1947), no. 62, 351--363.
\bibitem{yamagishi:95a} Masakazu Yamagishi, \textit{On the number of Galois p-extensions of a local field}, Proc. Amer.
Math. Soc. \textbf{123} (1995), no. 8, 2373--2380.

\end{thebibliography}
\end{document}